%% =====================================================================
%% Proposta de Trabalho Final - EEE017 Confiabilidade de Sistemas
%% UFMG - Prof. Michel Bessani
%% \textit{Dataset}: Bruinsma et al. (2024) - Motor + Bomba Centrífuga
%% Template: IEEEtran (conference)
%% Compilar com: pdflatex (2x) ou utilizar Overleaf
%% =====================================================================
\documentclass[conference]{IEEEtran}
\IEEEoverridecommandlockouts

% --- Pacotes essenciais ---------------------------------------------------
\usepackage[utf8]{inputenc}
\usepackage[T1]{fontenc}
\usepackage[brazil]{babel}
\usepackage{amsmath,amssymb,amsfonts}
\usepackage{graphicx}
\usepackage{cite}
\usepackage{url}
\usepackage{textcomp}
\usepackage{xcolor}
\usepackage{booktabs}
\usepackage{hyperref}
\hypersetup{
    colorlinks=true,
    linkcolor=black,
    citecolor=black,
    urlcolor=blue
}

\def\BibTeX{{\rm B\kern-.05em{\sc i\kern-.025em b}\kern-.08em
    T\kern-.1667em\lower.7ex\hbox{E}\kern-.125emX}}

\begin{document}

% =========================================================================
% TÍTULO E AUTORES
% =========================================================================
\title{Análise de Confiabilidade de Bombas Centrífugas Industriais
       Acionadas por Motor Elétrico com Base em
       Dados Experimentais de Sensores de Vibração}

\author{
\IEEEauthorblockN{Felipe Costa Lopes}
\IEEEauthorblockA{\textit{UFMG} \\
Matrícula: 2018019648 \\
Belo Horizonte, MG, Brasil}
\and
\IEEEauthorblockN{Isabella Beatriz de Souza Gomes}
\IEEEauthorblockA{\textit{UFMG} \\
Matrícula: 2022421587 \\
Belo Horizonte, MG, Brasil}
\and
\IEEEauthorblockN{Stéphanie Pereira Barbosa}
\IEEEauthorblockA{\textit{UFMG} \\
Matrícula: 2021088965 \\
Belo Horizonte, MG, Brasil}
}

\maketitle

% =========================================================================
% RESUMO
% =========================================================================
\begin{abstract}
Este trabalho propõe a aplicação de métodos clássicos de Engenharia de
Confiabilidade a um conjunto de dados experimental público de bombas
centrífugas industriais acionadas por motores de indução, coletado pela
Marinha Real dos Países Baixos~\cite{bruinsma2024}. O \textit{dataset} original contém
medições de vibração e corrente elétrica. Para os modos de falha estudados neste trabalho (defeitos de rolamento, desbalanceamentos e desalinhamento), utiliza-se exclusivamente o sinal de vibração. A partir desses dados, serão construídas séries de
\textit{tempos até falha} (TTF) por tipo de modo de falha, sobre as
quais serão ajustadas as distribuições Exponencial, Weibull e
Log-normal pelo método de ranqueamento com \textit{median ranks} de
Benard e por máxima verossimilhança (\textit{MLE}), com tratamento adequado de
dados censurados à direita. Intervalos de confiança para os parâmetros
serão estimados via \textit{bootstrap} não paramétrico. Com base na
distribuição ajustada, serão calculados os índices MTTF, taxa de falha
$\lambda(t)$, função de confiabilidade $R(t)$ e disponibilidade $A$.
Por fim, o conjunto motor-bomba será modelado por Diagrama de Blocos de
Confiabilidade (DBC) com os subsistemas \textit{motor} e \textit{bomba}
em série, avaliado via simulação de Monte Carlo. Cenários hipotéticos
de redundância ativa e \textit{standby} são incluídos como exercício de
projeto para quantificar o ganho de disponibilidade potencial. Espera-se
quantificar a influência de cada modo de falha na vida útil do
equipamento e fornecer recomendações de estratégia de manutenção.
\end{abstract}

\begin{IEEEkeywords}
Confiabilidade, bombas centrífugas, motor de indução, distribuição
Weibull, \textit{bootstrap}, Monte Carlo, diagrama de blocos,
manutenção preditiva.
\end{IEEEkeywords}

% =========================================================================
% I. INTRODUÇÃO
% =========================================================================
\section{Introdução}

A Engenharia de Confiabilidade estuda a probabilidade de um sistema
executar sua função prevista, sob condições especificadas, durante um
determinado período de tempo~\cite{oconnor2012}. Em sistemas industriais,
falhas não planejadas geram paradas de produção, custos elevados de
reparo e, em casos extremos, riscos à segurança. Por isso, a estimativa
de índices como o Tempo Médio até Falha (MTTF), a taxa de falha
$\lambda(t)$ e a disponibilidade $A$ é fundamental para a tomada de
decisão em manutenção.

Bombas centrífugas acionadas por motores de indução são equipamentos
críticos em praticamente todos os setores industriais --- da manufatura
ao naval~\cite{bruinsma2024}. Esses conjuntos motor-bomba estão sujeitos
a uma ampla variedade de modos de falha: defeitos de rolamento
(identificados pelas frequências BPFO e BPFI), desalinhamento, 
desbalanceamento, cavitação, dano no impelidor e degradação de
acoplamento~\cite{bruinsma2024}. A pergunta central do ponto de vista
da confiabilidade é: \textit{dado um modo de falha, qual é a vida
esperada do equipamento e como a configuração do sistema afeta a sua
disponibilidade?}

Este trabalho propõe responder a essa pergunta aplicando os métodos
ensinados na disciplina EEE017 --- ajuste de distribuições de vida,
\textit{bootstrap}, cálculo de índices e simulação de Monte Carlo --- a um
\textit{dataset} experimental público real de alta qualidade, descrito na
próxima seção.

% =========================================================================
% II. DESCRIÇÃO DO SISTEMA
% =========================================================================
\section{Descrição do Sistema}

\subsection{Dataset Utilizado}

O conjunto de dados utilizado é o \textit{Motor Current and Vibration
Monitoring \textit{Dataset} for Various Faults in an E-motor-driven Centrifugal
Pump}, publicado por Bruinsma \textit{et al.}~\cite{bruinsma2024} e
disponível gratuitamente no repositório 4TU.ResearchData (DOI:
\texttt{10.4121/2b61183e-c14f-4131-829b-cc4822c369d0}). O \textit{dataset} é
resultado de uma campanha experimental conduzida no \textit{Techport Fieldlab}
da Tata Steel, nos Países Baixos, pela Marinha Real Holandesa em
parceria com a Universidade de Twente.

O sistema físico consiste em duas bancadas completas de motor-bomba
centrífuga (Motor~2 e Motor~4), cada uma composta por:

\begin{itemize}
    \item Motor de indução trifásico (11\,kW a 1470\,RPM para o
    Motor~2; 22\,kW a 2950\,RPM para o Motor~4), acionado por
    inversor de frequência (\textit{Variable Frequency Drive} -- \textit{VFD});
    \item Bomba centrífuga NK80 acoplada ao motor por acoplamento de
    garras com blocos de borracha;
    \item Sistema de circulação de água (reservatório de 60\,m$^3$,
    válvulas de sucção e descarga controláveis).
\end{itemize}

\subsection{Falhas Simuladas}

O \textit{dataset} original registra o comportamento do sistema em condição saudável
(\textit{baseline}) e sob 11~tipos de falhas simuladas~\cite{bruinsma2024}.
Para o escopo e execução deste trabalho, foram selecionados e analisados
apenas 4~modos de falha específicos:

\begin{itemize}
    \item Defeito de rolamento (Motor~2): dano na pista
    interna (BPFI), em 3 severidades;
    \item Desbalanceamento no motor (Motor~4): em 6 severidades;
    \item Desbalanceamento na bomba (Motor~4): em 3 severidades;
    \item Desalinhamento paralelo (Motor~4): em 4 severidades.
\end{itemize}

\textit{Nota:} Os outros modos de falha presentes no \textit{dataset} (como 
dano em pista externa BPFO, desalinhamento angular, cavitação, dano no 
impelidor, pé solto, entre outros) foram avaliados na fase de planejamento, 
mas descartados deste estudo por não apresentarem assinatura de degradação 
suficientemente separável da condição base nos canais instrumentados ou 
por restrição de escopo do projeto.

\subsection{Sensores e Variáveis}

Duas técnicas de monitoramento por condição foram empregadas
simultaneamente~\cite{bruinsma2024}:

\begin{itemize}
    \item Vibração: 5~acelerômetros monoaxiais (Wilcoxon
    786B-10, 100\,mV/g), taxa de amostragem de 20\,kHz, cada
    registro com 12\,s de duração. Os canais cobrem os mancais
    do motor e da bomba em diferentes direções (horizontal,
    vertical e axial);
    \item Corrente e tensão: 3~fases medidas por pinças
    amperimétricas (CR3110) e taps de tensão (Wago 855), taxa de
    20\,kHz, registros de 15\,s por medição.
\end{itemize}

Esses dados têm rastreabilidade física completa: cada arquivo CSV é
identificado pelo motor, velocidade, tipo de falha e severidade, com
rotulagem quase perfeita (\textit{near perfect labelling})~\cite{bruinsma2024}. As variáveis medem os mesmos
fenômenos físicos presentes em motobombas industriais reais, tornando
o \textit{dataset} adequado para análise de confiabilidade com base em condição.

% =========================================================================
% III. REVISÃO DA LITERATURA
% =========================================================================
\section{Revisão da Literatura}

A análise estatística de dados de vida de equipamentos --- tempos até
falha (TTF) e dados censurados --- é tratada de forma abrangente em
O'Connor e Kleyner~\cite{oconnor2012} e Colosimo e
Giolo~\cite{colosimo2006}. As distribuições Exponencial (taxa de falha
constante), Weibull (taxa variável) e Log-normal (desgaste cumulativo)
são as mais utilizadas em confiabilidade, e seus ajustes por
ranqueamento com \textit{median ranks} de Benard e por \textit{MLE} com suporte
a censura à direita são procedimentos estabelecidos~\cite{oconnor2012}.

A distribuição Weibull, em particular, é versátil porque o parâmetro de
forma $\beta$ revela o estágio do ciclo de vida do equipamento: $\beta
< 1$ indica falhas prematuras (mortalidade infantil), $\beta \approx 1$
indica taxa de falha constante (vida útil), e $\beta > 1$ indica
desgaste crescente~\cite{bessani2026,oconnor2012}. Isso permite
relacionar o modo de falha ao comportamento esperado da taxa de falha
e à posição na curva da banheira.

O método \textit{bootstrap} não paramétrico para construção de
intervalos de confiança em parâmetros estimados de distribuições de
vida é apresentado como ferramenta computacional robusta nas notas de
aula da disciplina~\cite{bessani2026} e fundamentado em Robert e
Casella~\cite{robert2010}.

A modelagem de sistemas compostos por múltiplos componentes é feita por
meio do Diagrama de Blocos de Confiabilidade (DBC), que permite calcular
a confiabilidade e a disponibilidade do sistema nas configurações série,
paralelo com redundância ativa e \textit{standby}~\cite{oconnor2012,
bessani2026}. A simulação de Monte Carlo complementa essa análise em
sistemas onde a expressão analítica é complexa~\cite{zio2013,bessani2026}.

No campo de motores e bombas, Bruinsma \textit{et al.}~\cite{bruinsma2024}
demonstram que vibração e corrente são as técnicas de monitoramento mais
aplicadas na detecção precoce de falhas em equipamentos rotativos, o que
motiva o uso desses sinais como base para a derivação de tempos de vida
e consequente análise de confiabilidade.

% =========================================================================
% IV. METODOLOGIA
% =========================================================================
\section{Metodologia}

A metodologia é dividida em quatro etapas, implementadas em Python com
base nos scripts de referência da disciplina.

\subsection{Construção das Séries de Tempo até Falha (TTF)}

O \textit{dataset} de Bruinsma \textit{et al.}~\cite{bruinsma2024} registra, para
cada par (tipo de falha, severidade, velocidade), múltiplas medições de
12\,s de vibração (registros independentes por coluna do CSV). O procedimento para obtenção do
TTF é o seguinte:

\begin{enumerate}
    \item Para cada modo de falha $m$ e severidade $s$, computar o
    indicador de degradação: RMS do sinal de vibração em uma banda
    de frequência selecionada por física e validada por correlação de
    Spearman com a severidade (banda alta 1--9\,kHz para rolamento;
    banda baixa 10--300\,Hz para desbalanceamento e desalinhamento);
    \item Modelar a degradação como crescimento exponencial
    $D(\tau)=D_0 e^{\rho\tau}$, calibrando a taxa média $\bar\rho$ dos
    dados e simulando variabilidade unidade-a-unidade via taxa aleatória
    $\rho=\bar\rho\cdot e^{\mathcal{N}(0,\,c)}$; o equipamento é
    declarado em falha quando $D$ cruza o limiar
    $L=\mu_0+0{,}5\,(D_{max}-\mu_0)$;
    \item Registrar o TTF como o tempo decorrido (em minutos de
    operação acumulada) desde o início da medição em condição
    degradada até o cruzamento do limiar;
    \item Tratar como dado censurado à direita todo caso em
    que o limiar não foi ultrapassado ao fim do experimento, seguindo
    a metodologia apresentada na disciplina~\cite{bessani2026,colosimo2006}.
\end{enumerate}

Serão analisados os modos de falha com indicador de degradação
monotônico verificado: rolamento (BPFI), desbalanceamento no motor,
desbalanceamento na bomba e desalinhamento paralelo. O modo BPFO e o
desalinhamento angular foram avaliados, mas descartados por não
apresentarem assinatura separável da condição saudável nos canais disponíveis.

\subsection{Ajuste de Distribuições}

Sobre as amostras de TTF de cada modo de falha, serão ajustadas as
distribuições Exponencial, Weibull (2 parâmetros) e Log-normal por dois
métodos complementares, conforme apresentado na disciplina:

\paragraph{Ranqueamento com \textit{Median Ranks} de Benard}
Os TTFs observados são ordenados de forma crescente e seus ranques
medianos são calculados pela fórmula~\cite{oconnor2012,bessani2026}:
\begin{equation}
    r_{\text{med}}(i) = \frac{i - 0{,}3}{N + 0{,}4}
\end{equation}
onde $i$ é a posição do $i$-ésimo TTF na amostra ordenada e $N$ é o
total de observações (incluindo censuradas). Para dados censurados, os
ranques são corrigidos conforme o procedimento de Johnson apresentado
em~\cite{oconnor2012}. Os parâmetros são estimados por regressão linear
no gráfico de probabilidade de cada distribuição.

\paragraph{Máxima Verossimilhança (\textit{MLE})}
A estimativa por \textit{MLE} é realizada com suporte a censura à direita,
utilizando a biblioteca \texttt{reliability} do Python, que implementa
o método apresentado na disciplina. A comparação entre distribuições
será feita pelo critério AIC e pelos testes de aderência
Kolmogorov--Smirnov e Anderson--Darling.

\subsection{Intervalos de Confiança por \textit{Bootstrap}}

Para cada parâmetro estimado ($\beta$, $\eta$ na Weibull; $\mu$,
$\sigma$ na Log-normal; $\lambda$ na Exponencial), serão construídos
intervalos de confiança de 95\% por \textit{bootstrap} não paramétrico
com $R = 10^{4}$ reamostragens, seguindo o procedimento apresentado
nas notas de aula da disciplina e exemplificado no script
\texttt{6\_2\_IC\_Bootstrap.py}~\cite{bessani2026}:

\begin{enumerate}
    \item Reamostrar com reposição a amostra de TTFs (preservando
    o status de censura de cada observação);
    \item Reajustar a distribuição escolhida a cada reamostra;
    \item Calcular o percentil 2,5\% e 97,5\% das distribuições dos
    estimadores ao longo das $R$ repetições.
\end{enumerate}

\subsection{Cálculo dos Índices de Confiabilidade}

Com a distribuição ajustada $F(t)$, serão calculados os índices de
confiabilidade para cada modo de falha estudado~\cite{bessani2026,oconnor2012}:

\begin{itemize}
    \item MTTF: tempo médio até a falha,
    $\text{MTTF} = \int_0^\infty R(t)\,dt$, onde
    $R(t) = 1 - F(t)$ é a função de confiabilidade;
    \item Função de confiabilidade $R(t)$: probabilidade de
    sobrevivência até o instante $t$;
    \item Função de risco $h(t) = f(t)/R(t)$: taxa
    instantânea de falha, cujo comportamento indica o estágio na
    curva da banheira (crescente, constante ou decrescente);
    \item Disponibilidade $A = \text{MTTF}/(\text{MTTF} +
    \text{MTTR})$, com $\text{MTTR}=24$\,h (estimativa de engenharia
    para troca de rolamento/balanceamento, um turno de manutenção).
\end{itemize}

\subsection{Modelagem do Sistema e Simulação de Monte Carlo}

O conjunto motor-bomba será representado por um Diagrama de Blocos de
Confiabilidade (DBC) com dois subsistemas: o motor de indução
(sujeito a falhas elétricas e de rolamento do motor) e a
bomba centrífuga (sujeita a defeitos de impelidor, rolamento
da bomba e cavitação). Essa divisão é fiel à topologia física, na qual
a falha de qualquer um dos dois subsistemas interrompe o processo.

Serão avaliados três cenários, implementados em Python com base no
script \texttt{7\_2\_MC\_sistema\_dependencia.py} da
disciplina~\cite{bessani2026}:

\begin{itemize}
    \item Cenário 1 -- Série (configuração real): motor e
    bomba operam em série; qualquer falha para o sistema. Este é o
    arranjo físico efetivo da bancada;
    \item Cenário 2 -- Paralelo (redundância ativa)
    (hipotético): assume-se que um segundo conjunto motor-bomba
    idêntico está disponível; o processo mantém-se enquanto ao menos
    um dos dois conjuntos operar;
    \item Cenário 3 -- \textit{Standby} (hipotético): um conjunto
    opera e o segundo fica em espera, sendo ativado apenas em caso de
    falha do primeiro.
\end{itemize}

Os cenários~2 e~3 são exercícios de projeto, não descrevem a bancada
experimental, mas permitem quantificar o ganho de disponibilidade que
uma estratégia de redundância traria ao sistema real.

Para cada cenário, serão realizadas $N = 10^4$ iterações de Monte
Carlo, gerando TTFs por inversão da CDF ajustada:
\begin{equation}
    t = F^{-1}(U), \quad U \sim \mathcal{U}(0,1)
\end{equation}
e registrando se o sistema cumpre a missão no horizonte de tempo
estipulado. As expressões analíticas de disponibilidade para cada
configuração, apresentadas na disciplina~\cite{bessani2026}, serão usadas
como verificação dos resultados da simulação:
\begin{align}
    A_{\text{série}} &= \prod_{i} A_i \\
    A_{\text{paralelo}} &\approx 1 - \prod_{i}(1 - A_i) \\
    A_{\text{\textit{standby}}} &= \frac{\mu^2 + \mu\lambda}{\mu^2 + \mu\lambda + \lambda^2}
\end{align}

% =========================================================================
% V. RESULTADOS ESPERADOS
% =========================================================================
\section{Resultados Esperados}

Ao final do trabalho, espera-se entregar:

\begin{itemize}
    \item Tabela comparativa de MTTF, $\lambda$, $\beta$
    (Weibull) e disponibilidade $A$ para cada modo de falha analisado,
    permitindo identificar quais falhas degradam mais rapidamente o
    equipamento;
    \item Curvas $R(t)$ com bandas de confiança de 95\%
    (obtidas por \textit{bootstrap}), indicando, por meio do $\beta$
    estimado, em qual estágio da curva da banheira cada modo de falha
    se enquadra;
    \item Gráfico de probabilidade Weibull (ou da distribuição
    vencedora) para validação visual do ajuste e para comunicar os
    resultados de forma intuitiva para não especialistas;
    \item Comparação dos três cenários (série real, paralelo
    hipotético e \textit{standby} hipotético) em termos de
    $R_{\text{SYS}}(t)$ e disponibilidade do sistema, obtidos via
    simulação de Monte Carlo com verificação analítica;
    \item Recomendação de estratégia de manutenção: com base
    no $\beta$ estimado para cada modo de falha, será discutido se a
    manutenção preventiva baseada em tempo ou a manutenção baseada em
    condição é mais adequada.
\end{itemize}

Como contribuição metodológica, o trabalho demonstra uma rota
reprodutível para derivar indicadores de confiabilidade a partir de
dados experimentais de sensores, aplicando diretamente os métodos
da disciplina a um sistema físico real e bem documentado.

% =========================================================================
% REFERÊNCIAS
% =========================================================================
\begin{thebibliography}{99}

\bibitem{bruinsma2024}
S. Bruinsma, R.~D. Geertsma, R. Loendersloot, and T. Tinga,
``Motor current and vibration monitoring \textit{dataset} for various faults in
an E-motor-driven centrifugal pump,''
\textit{Data in Brief}, vol.~52, p.~109987, Feb. 2024,
doi: \url{10.1016/j.dib.2023.109987}.

\bibitem{oconnor2012}
P.~D.~T. O'Connor and A. Kleyner,
\textit{Practical Reliability Engineering}, 5th~ed.
Chichester: Wiley, 2012.

\bibitem{bessani2026}
M. Bessani,
``Notas de aula da disciplina EEE017 -- Confiabilidade de Sistemas,''
UFMG, Belo Horizonte, 2026.

\bibitem{colosimo2006}
E.~A. Colosimo and S.~R. Giolo,
\textit{Análise de Sobrevivência Aplicada}, 1ª~ed.
São Paulo: Blucher, 2006.

\bibitem{robert2010}
C.~P. Robert and G. Casella,
\textit{Introducing Monte Carlo Methods with R}.
New York: Springer, 2010.

\bibitem{zio2013}
E. Zio,
\textit{The Monte Carlo Simulation Method for System Reliability and
Risk Analysis}.
London: Springer, 2013.

\bibitem{birolini2007}
A. Birolini,
\textit{Reliability Engineering: Theory and Practice}, 5th~ed.
Berlin: Springer, 2007.

\bibitem{montgomery2013}
D.~C. Montgomery and G.~C. Runger,
\textit{Applied Statistics and Probability for Engineers}, 6th~ed.
Hoboken: Wiley, 2013.

\bibitem{pinciroli2023}
L. Pinciroli, P. Baraldi, and E. Zio,
``Maintenance optimization in industry 4.0,''
\textit{Reliability Engineering \& System Safety}, vol.~234,
p.~109204, 2023.

\end{thebibliography}

\end{document}